\subsection{Задания для лабораторной работы}
\subsubsection{Базовые миссии}

\begin{taskbox}[Полёт по окружности радиуса \texttt{r=0.5} на высоте \texttt{z=1.0}]
Требуется реализовать алгоритм, обеспечивающий устойчивый полёт по окружности:
\begin{itemize}
    \item выполнить \texttt{ap.push(Ev.MCE\_PREFLIGHT)} и \texttt{ap.push(Ev.MCE\_TAKEOFF)};
    \item запускать обновление точки через \texttt{Timer.new(0.1, fn)} после \texttt{Ev.TAKEOFF\_COMPLETE};
    \item вычислять \(x = r \cos(\theta)\), \(y = r \sin(\theta)\) с переводом градусов в радианы;
    \item при необходимости сопровождать курс вызовом \texttt{ap.updateYaw(angle * math.pi / 180)};
    \item использовать индикацию \texttt{Ledbar.new(29)} для этапов;
    \item по \texttt{Ev.SHOCK} — останов таймеров и \texttt{ap.push(Ev.ENGINES\_DISARM)}, по \texttt{Ev.LOW\_VOLTAGE2} — посадку.
\end{itemize}

\textbf{Ожидаемый результат:} квадрокоптер следует окружности радиуса \texttt{0.5} на высоте \texttt{1.0}, выполняет посадку, индикация отражает этапы.
\end{taskbox}

\begin{taskbox}[Полёт по многоугольнику из \texttt{N} точек окружности]
Требуется реализовать генерацию точек по равномерным углам \(a = 360 \cdot i / N\):
\begin{itemize}
    \item переводить углы в радианы при вычислении координат;
    \item стартовать маршрут по \texttt{Ev.TAKEOFF\_COMPLETE};
    \item переходить к следующей точке по \texttt{Ev.POINT\_REACHED};
    \item обеспечить безопасные реакции и индикацию как в окружности.
\end{itemize}

\textbf{Ожидаемый результат:} последовательный обход всех \texttt{N} точек и посадка после завершения маршрута.
\end{taskbox}

\begin{taskbox}[Траектория «восьмёрка»]
Требуется реализовать два состояния (\texttt{FIRST}, \texttt{SECOND}) для обхода двух окружностей с противоположным направлением:
\begin{itemize}
    \item сместить центры по оси X;
    \item выбрать шаг угла разного знака для фаз;
    \item завершить миссию возвратом в \texttt{(0,0,z)} и посадкой.
\end{itemize}

\textbf{Ожидаемый результат:} полёт по «восьмёрке» без рывков, завершение в центре и посадка.
\end{taskbox}

\begin{taskbox}[Спираль подъёма]
Требуется реализовать параметрическую спираль с приращениями \texttt{r = r + dr} и \texttt{z = z + dz} на каждом тике; по достижении \texttt{z\_max} остановить таймер, перейти в безопасную точку базы и выполнить посадку.

\textbf{Ожидаемый результат:} плавный подъём до \texttt{z\_max}, возврат к базе и посадка.
\end{taskbox}

\begin{taskbox}[Сопровождение курса \texttt{yaw}]
Требуется синхронизировать \texttt{ap.updateYaw} с параметром окружности, использовать радианы; ограничить шаг изменения угла разумным значением, чтобы исключить резкие скачки.

\textbf{Ожидаемый результат:} курс следует траектории без резких изменений, ориентация корректна.
\end{taskbox}

\begin{taskbox}[Запуск и останов миссии RC-тумблером]
Требуется считывать \texttt{Sensors.rc()}, анализировать \texttt{rc\_chans[8]}:
\begin{itemize}
    \item при включении тумблера выполнять предстарт и взлёт;
    \item при выключении — останавливаться, переходить в \texttt{(0,0,z)} и выполнять посадку;
    \item опрашивать RC через периодический \texttt{Timer.new(0.5, fn)}.
\end{itemize}

\textbf{Ожидаемый результат:} надёжный старт/останов миссии по положению тумблера с безопасным завершением.
\end{taskbox}

\begin{taskbox}[Аварийная реакция на \texttt{Ev.SHOCK}]
Требуется немедленно остановить все таймеры миссии, включить красный цвет на всех \texttt{0..28} индексах линейки и выполнить \texttt{ap.push(Ev.ENGINES\_DISARM)}.

\textbf{Ожидаемый результат:} мгновенное выключение двигателей и красная индикация аварии.
\end{taskbox}

\begin{taskbox}[Посадка при низком напряжении \texttt{Ev.LOW\_VOLTAGE2}]
Требуется остановить миссию, перейти в безопасную точку и выполнить \texttt{ap.push(Ev.MCE\_LANDING)}, включить красную индикацию низкого напряжения.

\textbf{Ожидаемый результат:} безопасная посадка при низком напряжении, отчёт индикацией.
\end{taskbox}

\begin{taskbox}[Обход квадрата по точкам, два цикла]
Требуется задать вершины \texttt{(-a/2,-a/2,z)}, \texttt{(a/2,-a/2,z)}, \texttt{(a/2,a/2,z)}, \texttt{(-a/2,a/2,z)}, выполнять переходы по \texttt{Ev.POINT\_REACHED} и совершить два полных обхода перед посадкой.

\textbf{Ожидаемый результат:} два полных обхода квадрата и посадка в конце маршрута.
\end{taskbox}

\begin{taskbox}[«Лоитер» по малой окружности \texttt{r=0.2}]
Требуется периодически обновлять точку окружности таймером, обеспечить безопасную остановку по событию или RC, сопровождать миссию световой индикацией этапов.

\textbf{Ожидаемый результат:} стабильное удержание по малой окружности и корректная остановка по команде.
\end{taskbox}
